\documentclass[12pt]{article}
\usepackage[T1]{fontenc}
\usepackage[utf8]{inputenc}
\usepackage{lmodern}
\usepackage{textcomp}
\usepackage{enumitem}
\usepackage{geometry}
\geometry{margin=1in}
\begin{document}
\begin{center}
Answer Key - Psychology - Undergraduate Level Assessment 3 \
\small (Answers are ordered exactly as in the question paper.)
\end{center}

\section*{Multiple Choice}
\begin{enumerate}[label=\arabic*.]
\item \textbf{Answer:} D
\\ \textit{Options:} A) when they are not on any form of medication; B) just before beginning antidepressant medication; C) during the first few days after starting antidepressant medication; D) as they begin to recover from their depression; E) during periods of remission from depression
\\ \textit{Note:} Correct option: D - as they begin to recover from their depression
\item \textbf{Answer:} A
\\ \textit{Options:} A) Random assignment, temporal precedence; B) Nonspuriousness, temporal precedence; C) Generalizability, causality; D) Temporal precedence, nonspuriousness; E) Temporal precedence, random selection
\\ \textit{Note:} Correct option: A - Random assignment, temporal precedence
\item \textbf{Answer:} E
\\ \textit{Options:} A) employee turnover but not job performance.; B) neither job performance nor health.; C) job performance but not health.; D) neither health nor employee turnover.; E) health but not job performance.
\\ \textit{Note:} Correct option: E - health but not job performance.
\item \textbf{Answer:} A
\\ \textit{Options:} A) family dynamics; B) career satisfaction; C) communication skills; D) emotional intelligence; E) insight
\\ \textit{Note:} Correct option: A - family dynamics
\item \textbf{Answer:} A
\\ \textit{Options:} A) To conduct psychological tests to diagnose the client; B) To prescribe medication for the client's mental health issues; C) To analyze the client's past traumas in every session; D) To find a cure for the client's issues; E) To judge and evaluate the client
\\ \textit{Note:} Correct option: A - To conduct psychological tests to diagnose the client
\end{enumerate}

\section*{True / False}
\begin{enumerate}[label=\arabic*.]
\setcounter{enumi}{5}
\item \textbf{Answer:} True
\\ \textit{Options:} A) True; B) False
\\ \textit{Note:} LLM-generated true statement based on MCQ.
\item \textbf{Answer:} True
\\ \textit{Options:} A) True; B) False
\\ \textit{Note:} LLM-generated true statement based on MCQ.
\item \textbf{Answer:} True
\\ \textit{Options:} A) True; B) False
\\ \textit{Note:} LLM-generated true statement based on MCQ.
\item \textbf{Answer:} False
\\ \textit{Options:} A) True; B) False
\\ \textit{Note:} LLM-generated false statement. Correct answer: '"No. it's not possible!"'.
\item \textbf{Answer:} False
\\ \textit{Options:} A) True; B) False
\\ \textit{Note:} LLM-generated false statement. Correct answer: 'no mode'.
\end{enumerate}

\section*{Long Form Response}
\begin{enumerate}[label=\arabic*.]
\setcounter{enumi}{10}
\item \textbf{Answer:} Cognitive dissonance is a learning theory that describes the process by which people adapt their behavior to match their intentions. Explain why this is the correct answer and provide supporting evidence. Reference key facts or concepts that support this choice.
\\ \textit{Note:} Long-form derived from MCQ; award credit for stating the correct idea and giving supporting reasoning.
\item \textbf{Answer:} Verbal materials are learned effectively if they are associated with strong emotional content. Explain why this is the correct answer and provide supporting evidence. Reference key facts or concepts that support this choice.
\\ \textit{Note:} Long-form derived from MCQ; award credit for stating the correct idea and giving supporting reasoning.
\end{enumerate}

\vfill
\noindent\textit{End of Answer Key}
\end{document}